\documentclass[11pt,a4paper]{article}
\usepackage[T1]{fontenc}
\usepackage[utf8]{inputenc}
\usepackage{amsmath,amssymb,amsthm}
\usepackage[margin=1in]{geometry}
\usepackage[colorlinks=true,linkcolor=blue,urlcolor=blue]{hyperref}
\theoremstyle{plain}
\newtheorem{theorem}{Theorem}
\newtheorem{lemma}[theorem]{Lemma}
\newtheorem{proposition}[theorem]{Proposition}
\newtheorem{corollary}[theorem]{Corollary}
\theoremstyle{definition}
\newtheorem{remark}[theorem]{Remark}

\title{The empirical recurrence for OEIS A186845, proved by transfer matrix}
\author{Adrian Perez Fontelles\\ \small Independent researcher}
\date{3 September 2026}
\begin{document}
\maketitle

\begin{abstract}
OEIS A186845 counts arrays over $\{0,\dots,1\}$ under a condition comparing a statistic of
every $2\times2$ subblock with those of its neighbours, and carries an empirical recurrence
of order $53$ contributed by R.~H.~Hardin. It is true, and it is decidable rather than
empirical. The count is a walk count on $S=512$ states and is $C$-finite; whether the
conjectured recurrence annihilates it is then settled by a finite exact computation.
\end{abstract}

\noindent\small 2020 Mathematics Subject Classification. 05A15, 05B45, 15B36.\normalsize

\section{The conjecture}

OEIS A186845 is ``Number of (n+1)X5 binary arrays with no 2X2 subblock diagonal sum less antidiagonal sum equal to any horizontal or vertical neighbor 2X2 subblock diagonal sum less antidiagonal sum.''. It has offset $1$ and begins
\[
512, 4194, 36988, 313412, 2679366, 22982034, 196096944, 1675514684,\ \dots
\]
The entry states, as a conjecture contributed by its author and never marked settled:
\begin{quote}
Empirical: a(n)=3*a(n-1)+22*a(n-2)+189*a(n-3)+260*a(n-4)+85*a(n-5)-1680*a(n-6)-4043*a(n-7)-4400*a(n-8)-13546*a(n-9)+13946*a(n-10)+56404*a(n-11)+230086*a(n-12)+252819*a(n-13)-187915*a(n-14)-799809*a(n-15)-2814268*a(n-16)-1329079*a(n-17)-1096497*a(n-18)+7541271*a(n-19)+13765860*a(n-20)+26386971*a(n-21)+5858979*a(n-22)-36979941*a(n-23)-92255551*a(n-24)-20494425*a(n-25)+20755236*a(n-26)+116599043*a(n-27)-16170165*a(n-28)+36077109*a(n-29)-62495789*a(n-30)+23873523*a(n-31)-36102678*a(n-32)+13081133*a(n-33)+835569*a(n-34)+230969*a(n-35)-719638*a(n-36)+4172847*a(n-37)-536371*a(n-38)+417666*a(n-39)-3123424*a(n-40)+1709226*a(n-41)-98800*a(n-42)+135240*a(n-43)-98051*a(n-44)-104636*a(n-45)+124283*a(n-46)-14637*a(n-47)-19679*a(n-48)+340*a(n-49)+788*a(n-50)+460*a(n-51)+152*a(n-52)-32*a(n-53) for n>55
\end{quote}
As of the ``Last modified'' line on the live entry (Oct 03 2025 03:44:37, revision 9) this is still
recorded as empirical, and nothing on the entry records it as proved.

\section{What is being counted}

Write $\sigma$ for its diagonal sum minus its antidiagonal sum, taken on a $2\times2$ subblock. The contiguous $2\times2$
subblocks of an $(n+1)\times5$ array over $\{0,\dots,1\}$ form a grid of $n$ rows
and $4$ columns, and $\sigma$ labels that grid. The entry
requires that no two subblocks that are neighbours horizontally or vertically carry the same value of $\sigma$.

\section{The count is a walk count}

The condition is on a PAIR of neighbouring subblocks, so nothing beyond the
rows themselves has to be remembered. Take as state the $2$-tuple of consecutive array rows
whose subblock row has just been completed: the horizontal pairs inside that row decide which
tuples are admissible at all, and the vertical pairs decide which tuple may follow which. A
step appends one array row. An array of $n+1$ rows is a walk of $n-1$ steps and
\[
a(n)\;=\;\iota^{\!\top}M^{\,n-1}\tau,\qquad\iota=\tau=(1,\dots,1)^{\!\top},
\]
on the $S=512$ admissible tuples.

The state does not compress. What a step needs of the past is the previous subblock row, and
for $5\ge2$ that row already determines the $2$ array rows it came from, so carrying the
rows costs nothing and carrying anything smaller loses information. Of the $2^{2\cdot
5}=1024$ tuples of $2$ rows, $S=512$ survive the condition inside a single subblock row;
the rest could not occur in a counted array at all.

Over the alphabet $\{0,\dots,1\}$ the statistic takes values in $[-2,2]$, so the
grid it labels carries at most $5$ different labels. At $n=1$ the array has $2$
rows and the grid is a single row of $4$
subblocks, with only the neighbours inside that row; the entry's first term is
$512$, which the model reproduces along with every later term, and that is the cheapest
check that the grid has been read with the right shape.

In particular $a$ is $C$-finite. The alignment of the walk index with the entry's own $n$ was
checked against its published terms rather than assumed.

\section{The proof}

\begin{lemma}\label{lem:crit}
Let $q(t)=t^{53}-\sum_i c_it^{53-i}$ be the characteristic polynomial of the
conjectured recurrence. The residual $a(n)-\sum_ic_ia(n-i)$ equals
$\iota^{\!\top}M^{\,j}q(M)\tau$ for the corresponding walk index $j$, so the recurrence
holds from some point on if and only if $\iota^{\!\top}M^{j}q(M)\tau=0$ for all large $j$,
and it is enough to examine $j<S$.
\end{lemma}

\begin{proof}
Factor $M^{j}$ out of $M^{j+53}-\sum_ic_iM^{j+53-i}$. For the bound, $M^{S}$
is an integer combination of $I,M,\dots,M^{S-1}$ by Cayley--Hamilton, so the numbers
$u_j=\iota^{\!\top}M^{j}q(M)\tau$ obey the monic recurrence given by the characteristic
polynomial of $M$; $S$ consecutive zeros force every later one, and the last nonzero $u_j$
pins the threshold exactly.
\end{proof}

\begin{theorem}
\[
a(n)\;=\;3\,a(n-1) + 22\,a(n-2) + 189\,a(n-3) + 260\,a(n-4) + 85\,a(n-5) - 1680\,a(n-6) - 4043\,a(n-7) - 4400\,a(n-8) - 13546\,a(n-9) + 13946\,a(n-10) + 56404\,a(n-11) + 230086\,a(n-12) + 252819\,a(n-13) - 187915\,a(n-14) - 799809\,a(n-15) - 2814268\,a(n-16) - 1329079\,a(n-17) - 1096497\,a(n-18) + 7541271\,a(n-19) + 13765860\,a(n-20) + 26386971\,a(n-21) + 5858979\,a(n-22) - 36979941\,a(n-23) - 92255551\,a(n-24) - 20494425\,a(n-25) + 20755236\,a(n-26) + 116599043\,a(n-27) - 16170165\,a(n-28) + 36077109\,a(n-29) - 62495789\,a(n-30) + 23873523\,a(n-31) - 36102678\,a(n-32) + 13081133\,a(n-33) + 835569\,a(n-34) + 230969\,a(n-35) - 719638\,a(n-36) + 4172847\,a(n-37) - 536371\,a(n-38) + 417666\,a(n-39) - 3123424\,a(n-40) + 1709226\,a(n-41) - 98800\,a(n-42) + 135240\,a(n-43) - 98051\,a(n-44) - 104636\,a(n-45) + 124283\,a(n-46) - 14637\,a(n-47) - 19679\,a(n-48) + 340\,a(n-49) + 788\,a(n-50) + 460\,a(n-51) + 152\,a(n-52) - 32\,a(n-53)
\]
for every $n>55$.
\end{theorem}

\begin{proof}
The vector $w=q(M)\tau$ was formed by $53$ matrix--vector products in exact integer
arithmetic and $\iota^{\!\top}M^{j}w$ evaluated for $j=0,1,\dots$, again exactly; those
integers vanish from the index corresponding to $n=55+1$ onwards and the run continued
until the working vector was identically zero, which settles every larger $j$ at once.
\end{proof}

\section{Verification}

Three checks, each independent of the algebra above.

First, the model reproduces all $17$ terms the entry publishes, exactly, in integer
arithmetic. This is what ties the model to the entry: the digraph is built by reading the
entry's English, and a misreading gives different counts.

Second, the conjectured recurrence was evaluated directly on those published terms, with no
matrices involved, and holds wherever the proved range and the data overlap.

Third, the annihilation test of Lemma~\ref{lem:crit} was carried out in exact integer
arithmetic on the whole vector rather than on a sampled prefix.

\begin{thebibliography}{9}
\bibitem{oeis} The OEIS Foundation, \emph{The On-Line Encyclopedia of Integer
Sequences}, \url{https://oeis.org}, 2026. Sequence A186845.
\bibitem{stanley} R.~P.~Stanley, \emph{Enumerative Combinatorics}, Volume 1, 2nd ed.,
Cambridge University Press, 2011. (Transfer-matrix method, Section 4.7.)
\bibitem{hu} J.~E.~Hopcroft, R.~Motwani and J.~D.~Ullman, \emph{Introduction to Automata
Theory, Languages, and Computation}, 3rd ed., Addison--Wesley, 2006.
\bibitem{horn} R.~A.~Horn and C.~R.~Johnson, \emph{Matrix Analysis}, 2nd ed., Cambridge
University Press, 2013.
\end{thebibliography}

\end{document}
